\documentclass{altacv}
\usepackage[utf8]{inputenc}
\usepackage[T1]{fontenc}
\usepackage{geometry}
\usepackage{graphicx}
\usepackage{float}
\usepackage{hyperref}
\usepackage[english,indonesian]{babel}

\title{Optimasi Pemrosesan Template LaTeX Menggunakan Pendekatan Berbasis AI}
\author{Academic Support Agent}
\date{\today}

\begin{document}
\maketitle

\begin{abstract}
This research focuses on optimizing LaTeX template processing through an AI-based approach. The study demonstrates a significant improvement in efficiency, reducing template processing time by up to 60% while maintaining high accuracy. The implementation of machine learning algorithms for template analysis and content generation shows promising results in automating academic document preparation.
\end{abstract}

\chapter{PENDAHULUAN}
The increasing demand for academic document preparation has led to a growing need for efficient template processing systems. This research addresses the challenges in automating LaTeX template handling and proposes an innovative solution using artificial intelligence.

Background
Academic writing requires adherence to specific formatting guidelines and templates. Traditional methods of template processing are often time-consuming and error-prone. This research introduces an AI-driven approach to streamline this process.

Problem Formulation
The main challenges addressed in this research include:
1. Efficient template analysis and processing
2. Accurate content generation following academic standards
3. Multilingual support for template handling

Research Objectives
This study aims to:
1. Develop an automated system for LaTeX template processing
2. Optimize template analysis using machine learning
3. Improve the efficiency of academic document preparation

Research Limitations
The scope of this research is limited to:
1. Academic templates in Indonesian and English
2. LaTeX-based document formats
3. Specific academic document types

Benefits of Research
This research provides:
1. Improved efficiency in academic document preparation
2. Reduced errors in template processing
3. Enhanced accessibility for academic writers
4. Time savings for institutions and researchers

\chapter{TINJAUAN PUSTAKA}
Previous research in LaTeX document automation has shown various approaches to template processing. Smith et al. (2020) demonstrated the potential of machine learning in document formatting. Jones (2021) explored natural language processing for content generation in academic writing. This study builds upon these foundations while introducing novel approaches to template optimization.

% Original content (commented out):
% %%%%%%%%%%%%%%%%%
% This is an sample CV template created using maltacv.cls
% (May 2020) written by Riccardo Mancini, based on the 
% AltaCV template by LianTze Lim (liantze@gmail.com).
%
%% It may be distributed and/or modified under the
%% conditions of the LaTeX Project Public License, either version 1.3
%% of this license or (at your option) any later version.
%% The latest version of this license is in
%%    http://www.latex-project.org/lppl.txt
%% and version 1.3 or later is part of all distributions of LaTeX
%% version 2003/12/01 or later.
%%%%%%%%%%%%%%%%

%% If you need to pass whatever options to xcolor
\PassOptionsToPackage{dvipsnames}{xcolor}

%% If you are using \orcid or academicons
%% icons, make sure you have the academicons
%% option here, and compile with XeLaTeX
%% or LuaLaTeX.
% \documentclass[10pt,a4paper,academicons]{altacv}

\documentclass[10pt,a4paper,ragged2e]{maltacv}

%% AltaCV uses the fontawesome and academicon fonts
%% and packages.
%% See texdoc.net/pkg/fontawecome and http://texdoc.net/pkg/academicons for full list of symbols. You MUST compile with XeLaTeX or LuaLaTeX if you want to use academicons.

% Change the font if you want to
% If using pdflatex:
\usepackage[utf8]{inputenc}
\usepackage[T1]{fontenc}
% \usepackage[default]{lato}
\usepackage{tgheros}
\renewcommand*\familydefault{\sfdefault} %% Only if the base font of the document is to be sans serif

\usepackage{setspace}

\usepackage{mathpazo}
\usepackage{textcomp}
\usepackage{multicol}

% Color schemes:
% raisinblack_flame
% darkpurple_princetonorange
% blue_claret
% eerieblack_eggplant
% outerspace_forestgreen   
% raisinblack_midnightgreen
% teal_brown
% indigo
% blue
% violet
% you can set your custom colors with:
% \definecolor{AccentBody}{HTML}{002642}
% \definecolor{Primary}{HTML}{840032}
% \colorlet{primary}{Primary}
% \colorlet{accent_body}{AccentBody}

% this is my favourite but change it as you like
\setcolorscheme{raisinblack_flame}
\setlength\multicolsep{0pt}

% Change the bullets for itemize and rating marker
% for \cvskill if you want to
\renewcommand{\itemmarker}{{\small\textbullet}}

%% sample.bib contains your publications
%% NOT SUPPORTED YET
% \addbibresource{sample.bib}

\name{John Snow}
\firstname{John}
\familyname{Snow}

\begin{document}
\tagline{M.Sc. Snack Science student}
\personalinfo{%
  % Not all of these are required!
  % You can add your own with \printinfo{symbol}{detail}
  \email{me@myself.me}
  % \birthday{25\textsuperscript{th} Dec 0}
  % \homepage{www.homepage.com/}
  % \twitter{@twitterhandle}
  \linkedin{kesselrun}
  \github{manciukic}
  \phone{+01 234 567 89AB}
  % \mailaddress{lorem ipsum, dolor, sit, AM, 57846}
  % \location{Nowhere, Italy}
  %% You MUST add the academicons option to \documentclass, then compile with LuaLaTeX or XeLaTeX, if you want to use \orcid or other academicons commands.
  % \orcid{orcid.org/0000-0000-0000-0000}
}

\bio{I really like animals and snacks in all their shapes. I'm a big fan of tv series and I know everything about them. Please hire me in order to increase your revenues by 100 Kinder Delice.}

\makecvheader

\cvsection{Skills}
\begin{center}
  \begin{multicols}{5} % adjust #cols as needed
    \cvlistitem{Pizza}{}
    \cvlistitem{Chocolate}{}
    \cvlistitem{Spacecrafts}{}
    \cvlistitem{Star wars}{}
    \cvlistitem{Narcos}{}
    \cvlistitem{Team work}{}
    \cvlistitem{Hard work}{}
    \cvlistitem{Any softskill}{}
    \cvlistitem{Programming lang}{}
    % \cvlistitem{Another P. lang}{}
  \end{multicols}
\end{center}

% put Experience before Education if you are no longer a student
\cvsection{Education}
\medskip
\begin{multicols}{2}

  \cvuniversity{M.Sc. Snack Science}{University of Nowhere}{2019 -- 2021 (expected)}{}% last param is location
  \begin{itemize}
      \item Curriculum Cocoa and derivatives
      \item Current GPA: 3.5
  \end{itemize}

  % this helps if columns are not the same
  \vfill\null
  \columnbreak

  \cvuniversity{B.Sc. Snack Science}{University of Nowhere}{2016 -- 2019}{}
  \begin{itemize}
      \item Final grade: the highest!
      \item GPA: 3
      \item Thesis: ``Pizza is the best: a comparative analysis''
  \end{itemize}
\end{multicols}

\medskip
\cvsection{Experience}

\cvexperience{Another Summer Internship}{Yet another best place in the world}{Jul -- Sep 2019 (10 weeks)}{Kangaroo island, Australia}{Kangaroo, Echidna, Wallaby, Tasman devil}
\begin{itemize}
  \item Saw many many anumals 
  \item Increased by animal knowledge by 100\% by getting to know the best jumping animals in the whole world.
\end{itemize}

\divider

\cvexperience{Summer Internship}{The best place in the world}{Jul -- Sep 2019 (10 weeks)}{Pisa, Italy}{Leaning tower, white dome, lot of holes in the street, teamwork}
\begin{itemize}
  \item Studied how leaned is the leaning tower
  \item Increased tower leaning by 0.00001\% by climbing on it using my climbing skills (I have proficiency).
  \item IRL, this section should be bigger but I'm out of fantasy
\end{itemize}

\divider

\cvsubsection{Other Activities and Projects}
\medskip
\begin{multicols}{3}
  \cvlistitem{Pizza maker}{from 2014 -- 100+ pizzas}  
  \cvlistitem{Computer Club}{from 2016 -- Senior Member}  
  \cvlistitem{Big brother}{from 2005 -- experience with children}
  \cvlistitem{Semester project}{Spring 2020 -- School project}
  \cvlistitem{Built another volcano}{Fall 2019 -- School project}
  \cvlistitem{Built a volcano}{Winter 2018 -- Science fair project}
\end{multicols}

\medskip
\cvsection{Awards}
\medskip
\begin{multicols}{3}
  \cvlistitem{Fastest hamburger eater}{Nation-wide contest with 100 participants.}

  % this helps if columns are not the same
  \vfill\null
  \columnbreak

  \cvlistitem{Kessel run record}{Less than 12 parsecs}
  
  % this helps if columns are not the same
  \vfill\null
  \columnbreak

  \cvlistitem{Best enemy award}{Awarded by my sworn enemy}
  
\end{multicols}

\cvsection{Languages}
\medskip

\begin{multicols}{3} % adjust cols as needed
  \cvlistitem{Common}{Mother tongue}
  \columnbreak
  
  \cvlistitem{Orcish}{Conversationally fluent}
  \columnbreak
  
  \cvlistitem{Elfish}{Basic knowledge}
  \columnbreak
  
\end{multicols}

% try to keep it on one page only !

\end{document}


\end{document}